\documentclass[11pt]{article}

\usepackage[margin=1in]{geometry}
\usepackage{amsmath}
\usepackage{amssymb}
\usepackage{booktabs}
\usepackage{hyperref}
\usepackage{enumitem}

\title{gPCD Geometry Dynamics}
\author{}
\date{\today}

\begin{document}
\maketitle

\section{Purpose}

The current Python dynamics model is an overlap-geometry collision model.  A
collision is not treated as an instantaneous impulse at first contact.  Instead,
the current overlap area determines where the pair is along a compression or
rebound path.

The old overlap-momentum response in \texttt{Dynamics.py} is no longer the
active collision response.  \texttt{Dynamics.py} remains as support code for
contact geometry, wall geometry, lifecycle helpers, and diagnostic overlap
momentum.  The active model is coordinated by \texttt{Base.py} and uses
\texttt{GeoDynamics.py} for the pair-level geometry calculation.

\section{Rules}

The model follows these project rules:

\begin{enumerate}[leftmargin=*]
    \item The model must remain an overlap-geometry model.
    \item Capture files are written as \texttt{CapXXXXXX.rpt} under
    \texttt{C:\textbackslash\_DJ\textbackslash gPCDData\textbackslash examples}.
    \item The model must remain parallel-friendly.
    \item Particle velocities in the cfg represent incoming velocity before
    turnaround; they define the particles' initial momentum.
    \item If a particle starts overlapped at time zero, the state is assumed to
    be post-turnaround and rebounding.
    \item \texttt{zero\_velocity\_overlap\_fraction} is good for designed
    natural collisions, but starting collisions need a calculated pair-level
    \(A_0\) because we do not know how deep into collision history they already
    are.
\end{enumerate}

\section{Contact Geometry}

For two particles, \texttt{Dynamics.particle\_contact()} returns:

\begin{itemize}[leftmargin=*]
    \item the unit normal \((n_x,n_y)\) from source center to target center,
    \item the current overlap area \(A\),
    \item the current center distance \(d\).
\end{itemize}

For equal-radius disks, the overlap area is:

\[
A(d) =
2r^2\cos^{-1}\left(\frac{d}{2r}\right)
- \frac{d}{2}\sqrt{4r^2-d^2}.
\]

Wall contacts use mirrored-circle geometry.  A particle at distance \(s\) from
a flat wall is treated as overlapping its mirror image at center distance
\(d=2s\).  The same circle-overlap area formula is then used.

\section{Zero-Velocity Overlap}

The zero-velocity overlap point, \(A_0\), is the deepest overlap reached by a
collision.  It is the turnaround point: compression ends and rebound begins.

For designed natural collisions, \(A_0\) comes from configuration:

\[
A_0 = f_0 A_{\max},
\]

where \(f_0=\texttt{zero\_velocity\_overlap\_fraction}\).

For starting overlaps, \(A_0\) is solved from incoming momentum and overlap
geometry.  If the solved \(A_0\) is shallower than the current overlap, the
model promotes \(A_0\) to the current overlap.  That keeps the starting state
physically possible.

\section{Velocity Profile}

Let \(A\) be the current overlap area and \(A_0\) be the zero-velocity overlap
area.  The compression fraction is:

\[
c = \mathrm{clamp}\left(\frac{A}{A_0}, 0, 1\right).
\]

The velocity fraction is:

\[
v_f = \sqrt{1-c}.
\]

This creates a parabolic collision process, analogous to the velocity profile
of a ball thrown upward in a vacuum.  At first contact, \(c=0\) and
\(v_f=1\).  At the zero-velocity point, \(c=1\) and \(v_f=0\).  During rebound,
the same geometry fraction maps velocity back from zero toward full rebound.

\section{Relative Velocity And Momentum}

For a particle pair, only relative velocity along the contact normal participates
in the collision response:

\[
v_n = (\vec{v}_t - \vec{v}_s)\cdot\vec{n}.
\]

The reduced mass is:

\[
\mu = \frac{m_s m_t}{m_s + m_t}.
\]

The incoming relative normal momentum is:

\[
p_n = \max(0, -\mu v_n).
\]

Tangential velocity is preserved.  The normal component is interpolated from the
starting velocity to the zero-velocity point during compression, and from the
zero-velocity point to full rebound during rebound.

\section{Contact State}

The Python model uses explicit contact state so a collision can remember its
history across frames.  This is essential because the rebound path cannot be
deduced from instantaneous geometry alone.  At the zero-velocity point, the
current velocity may be zero; without saved first-contact information, the model
would forget which way rebound should go.

\subsection{Where Contact State Lives}

Contact-state memory currently lives in \texttt{Base.py}, but particle-pair
state is now stored on the source particle:

\begin{itemize}[leftmargin=*]
    \item \texttt{particle["gcs"]}: source-owned particle-particle contact
    state slots.
    \item \texttt{Base.wall\_contact\_state}: particle-wall contact state.
\end{itemize}

\texttt{gcs} is initialized in \texttt{Base.reset()} for every live particle.
It is a fixed list with one available state slot per configured particle-contact
slot.  The current Python model uses \texttt{max\_contacts\_per\_particle}
slots, matching the existing \texttt{ccs} contact list.  Wall state remains in a
dictionary keyed by particle and wall flag.

\subsection{Particle-Pair Keys}

Particle-pair contact state is source-owned.  The key inside a particle is the
target particle index:

\[
\texttt{source particle gcs entry} \rightarrow \texttt{target\_index}.
\]

This means a physical pair can have two mirrored state records:

\[
P_i.\texttt{gcs} \rightarrow j
\qquad\text{and}\qquad
P_j.\texttt{gcs} \rightarrow i.
\]

The duplication is intentional.  It matches the future GPU ownership rule:
each source particle may write only its own state.  Normals are opposite in the
two records, but each source uses its own normal convention consistently.

Each state record stores:

\begin{itemize}[leftmargin=*]
    \item \texttt{active}: whether the slot is in use,
    \item \texttt{target\_index}: the target particle for this source-owned
    contact,
    \item \texttt{touched}: whether this state was observed in the current
    detection pass,
    \item \texttt{geo\_phase}: \texttt{compression} or \texttt{rebound},
    \item \texttt{first\_contact\_velocities}: starting velocity for both
    particles,
    \item \texttt{first\_contact\_normal},
    \item \texttt{first\_contact\_distance},
    \item \texttt{first\_contact\_overlap\_area},
    \item \texttt{geo\_zero\_velocity\_overlap\_area},
    \item \texttt{geo\_zero\_velocity\_center\_distance},
    \item diagnostic values such as relative normal momentum and overlap
    contact momentum.
\end{itemize}

When a \texttt{gcs} slot is not touched during a contact update, the contact has
exited.  If the state had reached rebound, the final rebound velocity is applied
before the slot is cleared.

\subsection{Wall Keys}

Wall contact state is keyed by:

\[
\texttt{wall\_key} = (\texttt{particle\_index}, \texttt{wall\_flag}).
\]

The wall flags are:

\begin{center}
\begin{tabular}{cl}
\toprule
Flag & Wall \\
\midrule
1 & left \\
2 & right \\
3 & bottom \\
4 & top \\
\bottomrule
\end{tabular}
\end{center}

Each wall state record stores:

\begin{itemize}[leftmargin=*]
    \item \texttt{geo\_phase},
    \item \texttt{first\_contact\_velocity},
    \item \texttt{first\_contact\_normal},
    \item \texttt{first\_contact\_distance},
    \item \texttt{first\_contact\_overlap\_area},
    \item \texttt{geo\_zero\_velocity\_overlap\_area},
    \item \texttt{geo\_zero\_velocity\_center\_distance}.
\end{itemize}

\subsection{Why This Memory Is Needed}

The geometry model needs three pieces of history:

\begin{enumerate}[leftmargin=*]
    \item the velocity at starting contact,
    \item the current phase, compression or rebound,
    \item the zero-velocity overlap geometry \(A_0\).
\end{enumerate}

The current position and velocity alone are not enough.  At turnaround the
normal velocity is zero, so a memoryless model can lock at the zero point or
misclassify the rebound as compression.  Contact state prevents that by keeping
the original first-contact velocity and the chosen \(A_0\).

\section{Step Flow}

The Python step flow in \texttt{Base.step()} is:

\begin{enumerate}[leftmargin=*]
    \item Begin the dynamics step.
    \item Detect current particle and wall contacts.
    \item Update particle-pair contact state.
    \item Update wall contact state.
    \item Apply geometry dynamics.
    \item Move particles.
    \item Detect contacts again at the new positions.
    \item Update wall contact state again.
    \item Clip rebound overlaps so contacts do not penetrate deeper than
    \(A_0\).
    \item End the dynamics step and record the report.
\end{enumerate}

\section{Starting Overlaps}

At reset, the model detects contacts before the first report.  Starting
particle-pair overlaps are initialized by
\texttt{initialize\_geo\_starting\_overlap\_zero\_state()}.  Starting wall
overlaps are initialized by
\texttt{initialize\_geo\_starting\_wall\_overlap\_zero\_state()}.

For starting overlaps:

\begin{itemize}[leftmargin=*]
    \item the phase is set to \texttt{rebound},
    \item \(A_0\) is solved from incoming momentum,
    \item if needed, \(A_0\) is promoted to the current overlap,
    \item the current velocity is set to the velocity that corresponds to the
    current overlap on the rebound path.
\end{itemize}

For starting wall overlaps, if the cfg velocity is already outward from the
wall, the incoming first-contact velocity is inferred by reflecting that
velocity across the wall normal.

\section{Exit Finalization}

When a contact exits, the state is removed.  If the contact had reached rebound,
the model finalizes the velocity to the full rebound velocity computed from the
stored first-contact velocity and normal.

For particle pairs, \texttt{finalize\_exited\_geo\_contact()} applies the final
pair rebound impulse.  For wall contacts,
\texttt{finalize\_exited\_geo\_wall\_contact()} reflects the stored incoming
wall velocity.

\section{Reporting}

Capture reports are compact text files.  They show:

\begin{itemize}[leftmargin=*]
    \item total momentum,
    \item kinetic energy,
    \item per-particle dynamics,
    \item particle-particle collision diagnostics,
    \item wall collision diagnostics.
\end{itemize}

These reports are intended to compare the live velocity against the expected
geometry-state velocity without reintroducing the old dynamics response.

\section{GPU Direction}

The Python \texttt{particle["gcs"]} contact-state slots are now the reference
shape for the GPU port.  The GPU should embed equivalent source-owned contact
state in the particle structure:

\begin{itemize}[leftmargin=*]
    \item one slot for each \texttt{ccs} particle contact,
    \item later, one slot for each \texttt{bcs} wall contact.
\end{itemize}

That keeps the GPU rule parallel-friendly: each invocation writes only its own
source particle state.  The state must store first-contact velocity, phase, and
\(A_0\), matching the Python dictionaries described above.

\end{document}
